% !TeX program = pdflatex
% =============================================================================
% Lektionsplan (Diskussionsgrundlage, NICHT das Arbeitsblatt selbst)
% Gewichtskraft & Federkraft -- UZH Praktikum I, Sara Romer Urban, KS im Lee, HS2026
% Klasse: 2e
% =============================================================================
\documentclass[11pt,a4paper]{article}

\usepackage[T1]{fontenc}
\usepackage[utf8]{inputenc}
\usepackage[ngerman]{babel}
\usepackage{lmodern}
\usepackage[a4paper,margin=2.2cm]{geometry}
\usepackage{array,booktabs,tabularx,xltabular}
\usepackage{enumitem}
\usepackage{xcolor}
\usepackage{amsmath}
\usepackage{siunitx}
\sisetup{per-mode=fraction,fraction-command=\frac}

\definecolor{titleblue}{HTML}{183A6B}
\definecolor{accentblue}{HTML}{2E6F9E}
\definecolor{groupteal}{HTML}{1B7F72}
\definecolor{darkgray}{HTML}{4A4A4A}

\setlength{\parindent}{0pt}
\setlength{\parskip}{0.5em}

\newcommand{\Methode}[1]{{\small\color{groupteal}\textbf{#1}}}

\usepackage{titlesec}
\titleformat{\section}{\Large\bfseries\color{titleblue}}{\thesection}{0.6em}{}
\titlespacing*{\section}{0pt}{1.2em}{0.5em}

\begin{document}
\sloppy

\begin{center}
{\Huge\bfseries\color{titleblue}Lektionsplan: Gewichtskraft \& Federkraft}\\[0.4em]
{\large Doppellektion (2$\times$45 Min.) -- Diskussionsgrundlage}
\end{center}
\vspace{0.5em}

\noindent\textbf{Klasse:} 2e \hfill \textbf{Lehrperson:} Loris Keller

\vspace{0.3em}
\noindent\textcolor{darkgray}{Dieses Dokument ist \textbf{keine} Reinschrift des Arbeitsblatts, sondern die
Planungsgrundlage dafür: Lernziele, Aufbau der Lektion und die gewählte
Didaktik-Methode pro Phase (aus \texttt{Methodensammlung\_Physikunterricht.pdf}).
Nach Absprache wird daraus das \texttt{.tex}-Arbeitsblatt lokal in VS Code
erstellt.}

% =========================================================================
\section*{1. Abgrenzung und Lernziele}

\textbf{Im Fokus:} Gewichtskraft ($F_G = m\cdot g$, ortsabhängig über den
Ortsfaktor $g$) und Federkraft nach dem Hookeschen Gesetz
($F = D\cdot\Delta x$, Federkonstante $D$ als Steigung im
$F$-$\Delta x$-Diagramm), verknüpft über das Prinzip des
Kräftegleichgewichts (Kalibrierung eines Kraftmessers).

\textbf{Bereits bekannt (nur Repetition, keine Neueinführung):} $F\sim a$ aus
dem Modellwagen-Experiment der letzten Stunde, das Grundprinzip eines
Kraftmessers (Feder + Skala in Newton, aus der Werkstatt Kraft), der
Ortsfaktor $g$ als Begriff.

\textbf{Bewusst nicht Teil dieser Einheit:} die Trägheits-/Impuls-Stationen
aus derselben Werkstatt (Münzen/Lineal, Crash-Test Bälle, Lokomotive/Magnete)
-- diese gehören zur späteren Wechselwirkung-Einheit (Newton 3) und werden
hier nur am Rande im Plenumsrückblick erwähnt. Ebenfalls nicht Teil:
Federspannenergie, Reihen-/Parallelschaltung von Federn, Luftwiderstand
(ein ursprünglich geplantes Hammer-und-Feder-Beispiel zur Mondlandung wurde
deshalb wieder vollständig gestrichen, siehe Abschnitt~3) sowie das
Federpendel (die Vertiefungsaufgabe \glqq Wiegen im Weltraum\grqq{} bleibt
deshalb bewusst beim Stichwort Trägheit, ohne den
Schwingungsmechanismus zu erklären).

\begin{itemize}[leftmargin=1.3cm,itemsep=0.3em]
  \item[LZ1] die Kraft-Relation $F\sim a$ auf ein neues Beispiel anwenden
  (Fussballschuss) und damit Kinematik (Beschleunigung) und Dynamik (Kraft)
  explizit verknüpfen.
  \item[LZ2] zwischen Trägheit (Masse) und Gewichtskraft eines Körpers
  unterscheiden und begründen, welche der beiden Grössen sich beim Wechsel
  des Ortes (z.\,B. von der Erde zum Mond) ändert und welche gleich bleibt.
  (Ergebnis bewusst nicht vorweggenommen, siehe Abschnitt 3.)
  \item[LZ3] die Gewichtskraft berechnen ($F_G=m\cdot g$) und mit
  unterschiedlichen Ortsfaktoren rechnen (Erde, Mond, ggf. unbekannter
  Planet), inkl. Wissen, dass $g$ gleichwertig in $\text{m/s}^2$ oder
  $\text{N/kg}$ angegeben werden kann.
  \item[LZ4] explizit zwischen der Längen\emph{änderung} $\Delta x$ einer
  Feder und ihrer absoluten Länge $x$ unterscheiden, insbesondere bei einer
  bereits vorgespannten Feder.
  \item[LZ5] aus eigenen Messdaten (Federkraftmesser-Experiment) den
  Zusammenhang zwischen Federkraft $F$ und Längenänderung $\Delta x$
  untersuchen, daraus die Federkonstante $D$ bestimmen und zwei Federn
  unterschiedlicher Härte miteinander vergleichen. (Ergebnis bewusst nicht
  vorweggenommen, siehe Abschnitt 3.)
  \item[LZ6] das Kalibrierungsprinzip eines Kraftmessers erklären
  (Gleichgewicht zwischen Feder- und Gewichtskraft) und es anwenden, um zu
  beschreiben, wie sich aus einer unkalibrierten Feder mit einfachen
  Mitteln ein funktionierender Kraftmesser bauen liesse.
  \item[LZ7] die Einheit der Federkonstante von N/cm in N/m umrechnen.
\end{itemize}

% =========================================================================
\section*{2. Aufbau der Doppellektion}

\begin{xltabular}{\textwidth}{@{}p{2.6cm}X p{1.2cm}@{}}
\toprule
\textbf{Phase} & \textbf{Inhalt \& Methode} & \textbf{Zeit} \\
\midrule
\endhead

Wiederholung &
\textbf{Plenum: Kraftformel \& Werkstatt-Rückblick.} Kurze Auffrischung von
$F\sim a$ (Modellwagen-Experiment) und der wichtigsten Beobachtungen aus der
Werkstatt Kraft (Kraftmesser-Grundprinzip; Trägheits-Stationen nur am Rande
erwähnt, nicht neu bearbeitet). \Methode{B1 Vergewisserungsphase} -- kurze
aktive Denkpause (allein/zu zweit), bevor im Plenum weitergemacht wird. &
6' \\[0.4em]

Übung &
\textbf{$F=ma$ am neuen Beispiel: Fussballschuss.} Ein Fussball hat vor dem
Schuss $v=0$, danach eine gewisse Geschwindigkeit -- welche Kraft muss der
Fuss ausgeübt haben (LZ1)? Explizite Verknüpfung Kinematik (Beschleunigung)
$\leftrightarrow$ Dynamik (Kraft). \Methode{C3 Think-Pair-Share} -- zuerst
allein überlegen, dann zu zweit austauschen, abschliessend im Plenum
teilen. &
10' \\[0.4em]

Aktivierung 1 &
\textbf{Steinstoss: Erde vs. Mond -- Trägheit.} Ein Astronaut schlägt sich
mit dem Fuss zuerst auf der Erde, dann gedanklich auf dem Mond am selben
Stein an: Tut es auf dem Mond mehr, weniger oder gleich weh? Kernpunkt: die
Trägheit des Steins ist auf beiden Himmelskörpern exakt gleich, weil
Trägheit eine Eigenschaft der Masse ist, nicht der Gravitation (PhysikLibre
Kap.~4.2: \glqq Die Masse $m$ \dots{} wird auch als träge Masse oder
Trägheit bezeichnet\grqq{}) (LZ2). \Methode{B1 Vergewisserungsphase} --
kurze Denkpause allein, danach Plenum; kein Video nötig, da nichts
Beobachtbares unterschiedlich ausfällt. &
4' \\[0.4em]

Aktivierung 2 &
\textbf{Golfball auf dem Mond -- separat anschliessend.} Neue, getrennte
Frage: Kommt ein Golfball auf dem Mond weiter als auf der Erde, und warum?
Hier geht es nicht mehr um Trägheit, sondern um die kleinere Gewichtskraft
(kleinerer Ortsfaktor $g$) (LZ2, Vorgriff LZ3). \Methode{B1
Vergewisserungsphase} -- Partnerarbeit statt Einzelarbeit (reine
Gedankenaufgabe, kein Experiment), danach Plenumsdiskussion, in der die
Lehrperson das historische Ergebnis (Shepards Golfschlag, Apollo~14, 1971)
bestätigt und die Erklärung über die kleinere Gewichtskraft sichert (kein
Video, da verfügbares Archivmaterial in der Bildqualität zu schlecht ist). &
4' \\[0.4em]

Theorie 1 &
\textbf{Gewichtskraft.} $F_G=m\cdot g$, Ortsfaktor $g$ für Erde und Mond,
Trägheit-Gewichtskraft-Unterschied aus beiden Aktivierungen festigen (LZ2,
LZ3). Statt einer Ortsfaktor-Tabelle zeigt eine Abbildung
(\texttt{Ortsfaktoren\_Planeten.jpg}, Quelle astronomiefans.de) die
Ortsfaktoren mehrerer Himmelskörper auf einen Blick; daran anschliessend
wird erklärt, dass $g$ gleichwertig in $\text{m/s}^2$ oder $\text{N/kg}$
angegeben werden kann (die Abbildung verwendet N/kg). Lehrervortrag,
geleitet, knapp gehalten. &
7' \\[0.4em]

Theorie 2 &
\textbf{Federkraft-Konzept, OHNE das Hookesche Gesetz vollständig
vorzutragen.} Kalibrierungsprinzip eines Kraftmessers: die Feder dehnt sich,
bis Feder- und Gewichtskraft im Gleichgewicht sind (LZ6). Vorspannung der
Feder und der Unterschied zwischen Längenänderung $\Delta x$ und Federlänge
$x$ werden eingeführt -- die \emph{quantitative} Beziehung $F=D\cdot\Delta x$
und $D$ selbst leiten die SuS in der Lernaufgabe aus eigenen Messdaten her
(siehe unten, analog zum Vorgehen bei Kondensatoren). Direkt im Anschluss
eine kurze individuelle Transferfrage (\glqq Vorgespannte Federn im
Alltag\grqq{}, z.\,B. Autofederung/Stossdämpfer): Nennen und begründen, wo
eine bereits vorgespannte Feder im Alltag vorkommt (LZ4). Lehrervortrag,
geleitet, knapp gehalten. &
9' \\[0.4em]

Übergang &
\textbf{Wechsel zur Partnerarbeit.} Umbau vom Plenum zu Partnerstationen:
Gruppen bilden, Stationsmaterial (Massstab, Federn, Massen, Laptop/Excel-
Vorlage) verteilen. Bewusst als eigene Zeile budgetiert, statt in der
Lernaufgabe unterzugehen -- dieser Wechsel kostet bei einer echten
Messapparatur (anders als bei einer Simulation) real Zeit. &
4' \\[0.4em]

Lernaufgabe &
\textbf{Federkraftmesser-Experiment (Partnerarbeit) -- SuS leiten
$F=D\cdot\Delta x$ und $D$ selbst her.} Das Aufgabenblatt öffnet mit einer
\glqq Ziel\grqq{}-Box (Untersuchungsfrage ohne Ergebnisvorwegnahme) und
einer Kräftediagramm-Aufgabe (Kräfte an der im Gleichgewicht hängenden
Masse selbst einzeichnen, als Brücke von Theorie~2 hierher). Danach
Stationsaufbau: Massstab mit zwei verschiebbaren Köpfen, oberer Kopf beim
Nullpunkt der ungespannten Feder gesetzt, Metallfuss zum Aufhängen, mehrere
Massen. Teil~1: Messreihen mit dem exakten Ortsfaktor $g=\SI{9.81}{\newton
\per\kilogram}$ aufnehmen, in die vorbereitete Excel-Vorlage eintragen
(Diagramm/Achsen bereits fertig), Steigung $D$ ablesen und die eigenen
N/cm-Werte zusätzlich von Hand in die SI-Einheit N/m umrechnen,
Rechenweg zeigen (LZ4, LZ5, LZ7 -- diese Umrechnung ist damit direkt in die
Lernaufgabe integriert, keine separate Übungsphase mehr nötig). Teil~2:
zwei Federn unterschiedlicher Härte vergleichen, dann eine
Konstruktionsaufgabe (aus der unkalibrierten, stärkeren Feder mit Stift,
bekannten Massen und einem Gehäuse selbst eine Federwaage bauen, LZ6).
Teil~3: mit dem eigenen $D$-Wert eine neue Auslenkung vorhersagen und am
Stationsaufbau überprüfen, dann als Synthese-Frage: Wie müsste die
Kalibrierung dieser Federwaage angepasst werden, wenn ein Astronaut sie auf
dem Mond benutzen soll (LZ2, LZ3, LZ6 verknüpft; beide Fälle -- Newton-Skala
bleibt unverändert korrekt, Kilogramm-Skala müsste umskaliert werden --
werden von den SuS anhand der beiden Gleichgewichtsgleichungen hergeleitet).
\Methode{D1 Lernaufgabe} -- die reale Messapparatur macht den linearen
Zusammenhang direkt erfahrbar; $F=D\cdot\Delta x$ entsteht hier bei den
Lernenden selbst, nicht im Theorie-Vortrag. Diese Phase ist das
didaktische Zentrum der Doppellektion und erhält deshalb bewusst den
grössten Zeitanteil, inklusive Auf-/Abbau der Federn, beider Messreihen und
der SI-Einheiten-Umrechnung. &
43' \\[0.4em]

Ausblick &
\textbf{Zusatzaufgaben.} Kurzer Hinweis auf das separate Dokument
\texttt{gewichtskraft-federkraft-zusatzaufgaben.tex} (Gewichtskraft auf
verschiedenen Himmelskörpern, Einheiten der Federkonstante an vorgegebenen
Werten, sowie drei an LEIFI angelehnte Vertiefungsaufgaben) für
Vor-/Nachbereitung. &
3' \\

\bottomrule
\end{xltabular}

\textcolor{darkgray}{\small Summe: 90 Min. Die frühere eigene Zeile
\glqq Übung: Einheitenumrechnung N/cm~$\to$~N/m\grqq{} entfällt, weil diese
Umrechnung neu direkt in Teil~1 der Lernaufgabe stattfindet (an den eigenen
Messdaten, mit gezeigtem Rechenweg); die frei gewordene Zeit ist in die
Lernaufgabe (neu 43') geflossen. Theorie~2 wurde um die kurze
Alltagstransfer-Frage zur Vorspannung auf 9' verlängert, der Ausblick im
Gegenzug auf 3' gekürzt, da \glqq Einheiten der Federkonstante\grqq{} und
\glqq Gewichtskraft berechnen\grqq{} nicht mehr im Arbeitsblatt selbst
erklärt werden müssen, sondern bereits als eigenständiges Zusatzdokument
vorliegen.}

\textcolor{darkgray}{\small Zusätzlich enthält das Arbeitsblatt eine
Vertiefungsaufgabe (\glqq Wiegen im Weltraum\grqq{}, eine einzelne Frage
zum Kalibrierungsprinzip in Schwerelosigkeit, ohne den
Federpendel-/Trägheitsmechanismus zu erklären), die bewusst nicht in den 90
Minuten budgetiert ist: als Zusatzaufgabe für schnelle SuS während der
Lernaufgabe oder als Hausaufgabe. Getrennt davon steht das eigenständige
Dokument \texttt{gewichtskraft-federkraft-zusatzaufgaben.tex} (5 Aufgaben,
28 Punkte, \texttt{\textbackslash noprintanswers}) für weitere
Vor-/Nachbereitung zur Verfügung, analog zu
\texttt{kondensatoren-zusatzaufgaben.tex}.}

% =========================================================================
\section*{3. Begründung der Methodenwahl}

Die Fussball-Aufgabe (C3 Think-Pair-Share) macht die Kinematik-Dynamik-
Verknüpfung an einem alltagsnahen Beispiel konkret, statt sie abstrakt zu
postulieren -- passend zur ausdrücklichen Vorgabe, diese Verbindung
herzustellen.

Die beiden Astronaut-Aktivierungen sind bewusst getrennt, weil sie
unterschiedliche Kernaussagen haben und sonst leicht vermischt würden.
Beide sind als reine Gedankenaufgaben angelegt (B1 Vergewisserungsphase),
nicht als Predict-Observe-Explain (M4): Es gibt in keiner der beiden
Aktivierungen etwas Beobachtbares, das ein Experiment oder ein Video
bräuchte. Bei Aktivierung~1 (Steinstoss) ist das offensichtlich, da die
Trägheit auf beiden Himmelskörpern identisch bleibt. Bei Aktivierung~2
(Golfball) wäre ein Video naheliegend gewesen (Shepards Golfschlag,
Apollo~14), aber das verfügbare Archivmaterial ist in der Bildqualität zu
schlecht, um im Unterricht tatsächlich etwas erkennen zu lassen. Die
Aktivierung bleibt deshalb ebenfalls eine reine Partnerdiskussion: Die SuS
begründen ihre Vermutung zu zweit, und die Lehrperson bestätigt das
historische Ergebnis erst in der anschliessenden Plenumsdiskussion, statt
es über ein schwaches Video zu belegen. Die Reihenfolge (zuerst Trägheit,
dann Gewichtskraft) bereitet ausserdem Theorie~1 vor, ohne die beiden
Konzepte in einer einzigen Aktivität zu vermengen.

Die Federkraftmesser-Lernaufgabe (D1) ist bewusst analog zur PhET-Lernaufgabe
bei Kondensatoren aufgebaut: Theorie 2 liefert nur das Konzept
(Kräftegleichgewicht, Vorspannung, $\Delta x$ vs. $x$), nicht die fertige
Formel. Die SuS leiten $F=D\cdot\Delta x$ und den Wert von $D$ selbst aus
ihren eigenen Messdaten her -- dasselbe Prinzip, das sich bei Kondensatoren
bewährt hat, hier auf ein reales Experiment statt eine Simulation
übertragen. \textbf{Entschieden:} Ja, konsequent bis in die Lernziele-Box
hinein. Weder die Formel $F=D\cdot\Delta x$ noch der Hinweis, dass $D$
einer Diagrammsteigung entspricht, taucht irgendwo im Arbeitsblatt vor der
Lernaufgabe auf, auch nicht in der Lernziele-Box: LZ4/LZ5 wurden dafür
bewusst kompetenzorientiert statt ergebnisoffen umformuliert (analog zu
LZ2, siehe Abschnitt~1).

Eine weitere Vertiefungsaufgabe wurde ergänzt, bewusst nicht als eigene
Zeile im Zeitplan budgetiert, sondern als Zusatzangebot: \glqq Wiegen im
Weltraum\grqq{} ersetzt eine frühere, rein erklärende Hinweisbox durch eine
echte Transferaufgabe. Die SuS begründen selbst, warum ein Federkraftmesser
in Schwerelosigkeit versagt, statt es nur zu lesen.

Ein ursprünglich geplantes Hammer-und-Feder-Beispiel (Apollo~15, David
Scott, 1971, zur Mondlandung-Simultaneität von Hammer und Feder) wurde
wieder gestrichen: Der Effekt ist nur unter Ausschluss von Luftwiderstand
korrekt erklärbar, und Luftwiderstand gehört nicht zum Thema dieser
Doppellektion.

% =========================================================================
\section*{4. Anschlussdokument: Aufgabenblatt}

Wie von Sara Romer gewünscht, entsteht zusätzlich ein separates
Aufgabenblatt: Es öffnet mit einer kurzen Kräftediagramm-Aufgabe (die
Kräfte an der im Gleichgewicht hängenden Masse selbst einzeichnen, als
Brücke von Theorie~2 zur Lernaufgabe) und enthält danach das komplette
Federkraftmesser-Stationsexperiment in drei Teilen: Messreihen aufnehmen,
die beiden Federn vergleichen, sowie eine Vorhersage mit dem eigenen
$D$-Wert und deren Überprüfung am Stationsaufbau.

% =========================================================================
\section*{5. Optionale Zusatzmaterialien (Selbstkontrolle, nicht Ersatz)}

Als optionale Ergänzung zum realen Federkraftmesser-Experiment, nicht als
Ersatz dafür: die PhET-Simulationen \glqq Hooke's Law\grqq{} und \glqq
Masses and Springs\grqq{} (Selbstkontrolle/Vertiefung, z.\,B. als
Hausaufgabe) sowie das phyphox-Experiment \glqq Spring\grqq{}, das
allerdings eher auf Schwingungsfrequenz zielt und damit näher an der
späteren Schwingungs-Einheit liegt als an der statischen Messung hier.

% =========================================================================
\section*{6. Offene Punkte zur Besprechung}

\begin{itemize}[leftmargin=1.2cm,itemsep=0.4em]
  \item Die Lernaufgabe wurde auf 38' (plus 4' eigene Übergangszeile)
  angehoben, damit der physische Stationsaufbau (Auf-/Abbau, zwei Federn,
  Excel-Eingabe) realistisch Platz hat -- dafür wurden Wiederholung,
  Fussball-Übung und beide Theorie-Blöcke gekürzt. Reicht das in der
  Praxis, oder braucht die Lernaufgabe noch mehr Zeit als bei einer reinen
  Simulation?
  \item Reicht eine einzige Werkstatt-Station (zwei Federn), oder sollen
  weitere Variationen (z. B. Gummiband als nichtlineares Gegenbeispiel)
  ergänzt werden?
\end{itemize}

\end{document}
